% =============================================================
%  ONLINE RESOURCE 1 — JAX Implementation Details
%  Supporting Information for Hackman & Zhu (2026)
%  "When Inverse Physics-Informed Neural Networks Become
%   Practically Identifiable: A Case Study in Synaptic Dopamine
%   Transport" — Medical & Biological Engineering & Computing
%   submission
% =============================================================
\documentclass[12pt]{article}

\usepackage[margin=1in]{geometry}
\usepackage[utf8]{inputenc}
\usepackage[T1]{fontenc}
\usepackage{lmodern}
\usepackage{amsmath, amssymb, mathtools}
\usepackage{hyperref}
\hypersetup{colorlinks=true, linkcolor=black, citecolor=black,
            urlcolor=blue, pdfborder={0 0 0}}
\usepackage{setspace}
\usepackage{lineno}
\usepackage{listings}
\linenumbers
\doublespacing
\setlength{\parindent}{0pt}
\setlength{\parskip}{0.7em}

\title{Online Resource 1: JAX Implementation Details \\
\large Supporting Information for ``When Inverse
Physics-Informed Neural Networks Become Practically
Identifiable: A Case Study in Synaptic Dopamine Transport''}

\author{Emmanuel Hackman \and Huiqing Zhu \\
School of Mathematics and Natural Sciences, \\
University of Southern Mississippi, Hattiesburg, MS, USA \\
Corresponding author: \texttt{Emmanuel.Hackman@usm.edu}}

\date{}

\begin{document}
\maketitle

We summarize the key JAX code blocks used to construct and train
the PINN models reported in the main text. The full reproducible
implementation, including plotting helpers, the analytical-solution
validation, and the explicit two-dimensional finite-difference
reference solver, is publicly available in the GitHub repository
cited in the Code Availability statement of the main text
(\url{https://github.com/Hachero98/dopamine-PINN-}). The
implementation uses Flax NNX for the network module, Optax for
the first-order Adam phase, and jaxopt for the second-order L-BFGS
phase. All collocation-point arrays are JAX-native
(\texttt{jnp.ndarray}) and the training step is JIT-compiled via
\texttt{nnx.jit}.

\section*{S1.1 Network, residual, and collocation sampling}

The MLP takes $(x, y, t) \in \mathbb{R}^3$ and returns a scalar
concentration. Spatial and temporal derivatives needed for the
PDE residual are obtained by composing \texttt{jax.grad} on the
network output; \texttt{jax.vmap} batches the per-point residual
across all $N_r$ collocation points without an explicit Python
loop.

\begin{verbatim}
import jax, jax.numpy as jnp
import numpy as np
from flax import nnx
import optax
from jaxopt import LBFGS

jax.config.update("jax_enable_x64", True)

D_TRUE, K_TRUE = 0.32, 0.05    # diffusion (mu m^2/ms), reuptake (1/ms)
L, T           = 5.0, 20.0     # side length (mu m), simulation (ms)
SIGMA, C0      = 0.5, 1.0      # release width (mu m), peak (mu M)

class MLP(nnx.Module):
    def __init__(self, layers, *, rngs):
        # Each Linear layer is its own named attribute (lin_0, lin_1, ...)
        # so Flax NNX treats them as named submodules without nnx.List.
        self.n_layers = len(layers) - 1
        for i in range(self.n_layers):
            setattr(self, f"lin_{i}",
                nnx.Linear(layers[i], layers[i + 1],
                           kernel_init=nnx.initializers.glorot_normal(),
                           rngs=rngs))
    def __call__(self, xyt):
        h = xyt
        for i in range(self.n_layers - 1):
            h = jnp.tanh(getattr(self, f"lin_{i}")(h))
        return getattr(self, f"lin_{self.n_layers - 1}")(h)

def C_at(model, x, y, t):
    return model(jnp.stack([x, y, t]))[0]

# PDE residual at one point: dC/dt - D Laplacian(C) + k C
def residual_one(model, x, y, t, D, k):
    dC_dt   = jax.grad(C_at, argnums=3)(model, x, y, t)
    d2C_dx2 = jax.grad(lambda u: jax.grad(C_at, argnums=1)
                                  (model, u, y, t))(x)
    d2C_dy2 = jax.grad(lambda v: jax.grad(C_at, argnums=2)
                                  (model, x, v, t))(y)
    return dC_dt - D * (d2C_dx2 + d2C_dy2) + k * C_at(model, x, y, t)
\end{verbatim}

\section*{S1.2 Forward problem}

The forward PINN treats $D$ and $k$ as known constants and learns
the concentration field $\hat{C}(x, y, t; \theta)$ that satisfies
the PDE residual together with the IC and Neumann BC losses. The
Adam phase runs in Optax; the L-BFGS phase runs through jaxopt on
the same NNX parameter pytree, exposed via \texttt{nnx.split} and
reassembled via \texttt{nnx.merge}.

\begin{verbatim}
rngs      = nnx.Rngs(1234)
model     = MLP([3, 64, 64, 64, 64, 1], rngs=rngs)
optimizer = nnx.Optimizer(model, optax.adam(1e-3), wrt=nnx.Param)

def forward_loss(model, pts):
    res = jax.vmap(residual_one,
                   in_axes=(None, 0, 0, 0, None, None))(
        model, pts["x_r"], pts["y_r"], pts["t_r"], D_TRUE, K_TRUE)
    L_r = jnp.mean(res ** 2)
    # IC and BC losses computed analogously (omitted for brevity)
    return 10.0 * L_r + L_ic + L_bc

@nnx.jit
def adam_step(model, optimizer, pts):
    loss, grads = nnx.value_and_grad(forward_loss)(model, pts)
    optimizer.update(model, grads)
    return loss

for it in range(20_000):
    adam_step(model, optimizer, pts)

# L-BFGS phase: NNX state becomes the parameter pytree for jaxopt
gdef, state = nnx.split(model)
def lbfgs_loss(params):
    return forward_loss(nnx.merge(gdef, params), pts)
result = LBFGS(fun=lbfgs_loss, maxiter=5_000, tol=1e-9).run(state)
model  = nnx.merge(gdef, result.params)
\end{verbatim}

\section*{S1.3 Inverse problem}

In JAX, the diffusion coefficient and the reuptake rate are
exposed as \texttt{nnx.Param} attributes of an
\texttt{InverseMLP} module, so they live in the same parameter
pytree as the network weights and are optimized jointly without
a dedicated ``external trainable variables'' interface.
Log-parametrization keeps $D$ and $k$ strictly positive.

\begin{verbatim}
class InverseMLP(nnx.Module):
    def __init__(self, layers, D0, k0, *, rngs):
        self.mlp   = MLP(layers, rngs=rngs)
        self.D_log = nnx.Param(jnp.asarray(np.log(D0)))
        self.k_log = nnx.Param(jnp.asarray(np.log(k0)))
    def __call__(self, xyt):
        return self.mlp(xyt)
    @property
    def D(self): return jnp.exp(self.D_log.value)
    @property
    def k(self): return jnp.exp(self.k_log.value)

def inverse_loss(model, pts, obs):
    D, k = model.D, model.k
    res = jax.vmap(residual_one,
                   in_axes=(None, 0, 0, 0, None, None))(
        model.mlp, pts["x_r"], pts["y_r"], pts["t_r"], D, k)
    L_r = jnp.mean(res ** 2)
    C_pred = jax.vmap(C_at, in_axes=(None, 0, 0, 0))(
        model.mlp, obs["x_d"], obs["y_d"], obs["t_d"])
    L_d = jnp.mean((C_pred - obs["C_d"]) ** 2)
    # IC and BC losses computed analogously (omitted for brevity)
    return 10.0 * L_r + L_ic + L_bc + L_d

# (obs_x, obs_y, obs_t, obs_C) drawn from FD reference at 2% noise.
# Deliberately offset initial guesses D_0 = 0.30, k_0 = 0.04.
model     = InverseMLP([3, 64, 64, 64, 64, 1],
                       D0=0.30, k0=0.04, rngs=nnx.Rngs(1234))
optimizer = nnx.Optimizer(model, optax.adam(1e-3), wrt=nnx.Param)

@nnx.jit
def step(model, optimizer, pts, obs):
    loss, grads = nnx.value_and_grad(inverse_loss)(model, pts, obs)
    optimizer.update(model, grads)
    return loss

history = []
for it in range(20_000):
    step(model, optimizer, pts, obs)
    if it % 500 == 0:
        history.append((it, float(model.D), float(model.k)))

# L-BFGS phase identical in shape to the forward problem
gdef, state = nnx.split(model)
def lbfgs_loss_inv(params):
    return inverse_loss(nnx.merge(gdef, params), pts, obs)
result = LBFGS(fun=lbfgs_loss_inv, maxiter=5_000, tol=1e-9).run(state)
model  = nnx.merge(gdef, result.params)
\end{verbatim}

After training, the physical parameters are recovered by reading
\texttt{model.D} and \texttt{model.k}, which apply
\texttt{jnp.exp} to the underlying log-parameterized variables.
The trajectory of $(D, k)$ logged every $500$ Adam iterations
is persisted to \texttt{variables.dat} for post-hoc visualization
(Fig 4 of the main text).

\end{document}
